% Included in model_v7.tex in place of the obsolete Task 2 section only.
\section{任务二：投资、调度与劳动力的动态优化}\label{sec:t2}

\paragraph{复核更正（2026-11-11）。}
原稿仅以累计面积为状态，将新增面积补水与全域稳态收益相配，遗漏了全期预算、
存量光伏现金流及完整投资分类。现改为按建植批次携带任务一的植被与有效储水状态，
在有限灌溉策略库内求解混合整数动态优化模型，
不再采用仅含累计面积的 Bellman 递推。原有线性补水响应和光伏面积占比阈值不再适用。
更正只涉及模型、参数和约束；计算结果另文报告。

\subsection{时间、状态和决策}
规划期 $t=0,\ldots,T-1$，$T$ 在题目建议的10--20年内选择。
一批生态治理单元用 $j=(s,p)$ 表示，其中 $s$ 为建植年，$p$ 为全寿命灌溉策略。
年新增面积与光伏容量分别按 $\Delta_A$ hm$^2$、$\Delta_P$ kW 离散。
$n_{sp}$、$z_s$ 为非负整数，面积 $a_{sp}=\Delta_A n_{sp}$，装机 $y_s=\Delta_Pz_s$。
同批次采用相同初态与气候驱动，面积只缩放水量、工时、费用和生态积分。
策略库须同时包含不补水、恒定、随林龄增加和随林龄减少的定额，避免先验固定“先多后少”。

年初建植并支付建设费，年内治理面积全部计养护，年末支付运维并更新植被。
光伏建设费在年初支付、年末投产，\textbf{下一年}才计发电与运维。
初始生态治理面积、光伏资产均为零；初始植被与储水是每批拟治理地的情景条件，
不是从优化中免费生成的实测水源。若初始储水需人工灌入，必须纳入该批供水账。
生态与光伏用地按不重叠处理，不凭空叠加板下植被收益。

\begin{align}
A_t&=\sum_{s\le t,p}a_{sp},\qquad a_t=\sum_p a_{tp},\\
P_t^{\rm inst}&=\sum_{s\le t}y_s,\qquad P_t^{\rm op}=\sum_{s<t}y_s.
\end{align}
系统状态为各在役批次 $(V_{sp,t},M_{sp,t})$、光伏分龄资产、累计用地、剩余预算。
原来仅有 $A_t$ 的状态不足以区分“新植”和“成熟”、“有水”和“缺水”地块。

\subsection{任务一接口与生态指标}
对每个策略的林龄 $h=t-s$ 年定额 $q_{p,h}$（mm/年），直接调用任务一复核后的水量模块：
\begin{equation}\label{eq:Ks}
(M_{sp,t+1},\omega_{sp,t})=\mathcal W(V_{sp,t},M_{sp,t},q_{p,t-s};\text{逐月气候与水分参数}),
\end{equation}
其中 $\omega$ 为实际/潜在蒸腾比，不是凭空拟合的补水线性响应；月内子步保留质量守恒。
\begin{equation}\label{eq:xstar2}
V_{sp,t+1}=V_{sp,t}+rV_{sp,t}(1-V_{sp,t}/K)\omega_{sp,t}
-dV_{sp,t}-kE_t(1-\eta V_{sp,t}/K)V_{sp,t}.
\end{equation}
基准继承任务一的 $k=0$ 情景；不能据此报告实测减蚀吨数。
生态积分取
\[
H=\sum_{t=0}^{T-1}\sum_{s\le t,p}a_{sp}\frac{V_{sp,t}}{K}\quad[\mathrm{hm^2\cdot year}].
\]
年内覆盖按任务一的年初状态保持，不能把年末生长量提前享受一年。
$H$ 是治理地覆盖达成率的面积时间积分，\textbf{不是因果新增生态效益、货币收益或实测输沙减量}。
未治理地反事实尚未标定，因此不能声称已扣除了自然恢复。

\subsection{资金分配、劳动和现金流}
将支出严格分成生态建设 $C_t^b$、全部运维补水 $C_t^o$、光伏收益设施建设 $C_t^p$。
设 $w_c$ 为光伏年清洗量（m$^3$/kW），$Q_t^e$ 为生态毛取水量。
\begin{align}
C_t^b&=c_ba_t,\qquad C_t^p=c_py_t,\\
Q_t&=Q_t^e+w_cP_t^{\rm op},\qquad
Q_t^e=10\sum_{s\le t,p}a_{sp}q_{p,t-s},\\
C_t^o&=c_mA_t+c_{pm}P_t^{\rm op}+c_wQ_t,\qquad
D_t=C_t^b+C_t^o+C_t^p,\\
(f_t^b,f_t^o,f_t^p)&=(C_t^b,C_t^o,C_t^p)/D_t\quad(D_t>0).
\end{align}
三类资金及比例由活动决策与上述等式联立确定；消去冗余资金变量不等于省略投资结构。
无支出年比例未定义，不伪造为三项同时零并声称比例和为一。
光伏运维归入运维，不重复归入收益设施建设。单价按含人工包价假设，劳动需求为
\[
L_t=\ell_ba_t+\ell_mA_t+\ell_py_t+\ell_{pm}P_t^{\rm op},
\]
单位为人日/年。$L_t$ 是与工作量相等的投入决策，约束可用劳动力；不再另加一遍工资。
工种、季节与雇工价格细化需要外部工效证据。

首个运行年的净上网比产出为 $e_0$，年衰减率为 $d_p$，净结算电价为 $p_e$：
\[
G_t=\sum_{s<t}y_se_0(1-d_p)^{t-s-1},\qquad R_t=p_eG_t.
\]
净上网口径已包含转换、积尘、限电等综合损失，不再重复乘同类折减。
没有按新增面积的一次性收入冒充全寿命发电收入。

\subsection{硬约束与量纲}
\begin{align}
&D_t\le B_t,\qquad \sum_tD_t\le B_{\rm tot},\label{eq:budget}\\
&Q_t\le Q_t^{\max},\qquad Q_{t,m}\le Q_{t,m}^{\max}\quad(m=1,\ldots,12),\label{eq:water}\\
&A_t+\alpha_pP_t^{\rm inst}\le A^{\max},\quad a_t\le a_t^{\max},
\quad P_t^{\rm inst}\le P^{\max},\label{eq:land}\\
&L_t\le L_t^{\max},\label{eq:labor}\\
&n_{sp},z_s\in\mathbb Z_{\ge0},\qquad p\in\mathcal P.\label{eq:discrete}
\end{align}
这里面积统一 hm$^2$，$1$ mm作用于 $1$ hm$^2$ 为 $10$ m$^3$；若面积改为千 hm$^2$，
体积才是 $Ia$ 万m$^3$，不得再除1000。取水涵盖\textbf{全部在役批次}，不仅是新增面积。
月取水按任务一的温度加权蒸散比例分配，清洗水均匀分月；这属于有限调度策略假设，
不能称逐月自由控制的最优解。降雨不计作人工取水。

预算是未贴现的全期\textbf{总支出}，收入不冲抵支出限额，也不循环再投资；
贴现只作用于目标。各年预算上限、水量上限及劳力上限必须同时满足。
可附加投资结构情景：
\[
\underline f_p\sum_tD_t\le\sum_tC_t^p\le\overline f_p\sum_tD_t.
\]
该全期比例约束与逐年三类比例不是同一个量。默认不强迫建设光伏。

\subsection{生态与成本的权衡及求解范围}
资本在年初发生，运维和收入在年末发生，$\delta=(1+i_d)^{-1}$，$i_d$ 为贴现率（区别于任务三退化阈值系数）。
定义净现成本
\begin{equation}\label{eq:npv}
N=\sum_{t=0}^{T-1}\left[\delta^t(C_t^b+C_t^p)+\delta^{t+1}(C_t^o-R_t)\right].
\end{equation}
不用任意价格把生态指标伪装成现金收入。采用 $\varepsilon$ 约束：
\begin{equation}\label{eq:value}
\min N\quad\text{s.t. }H\ge H_{\rm req},\quad\text{式}\ \eqref{eq:budget}\text{--}\eqref{eq:discrete}.
\end{equation}
先求相同资源约束下的最大生态参考，再以一组 $H_{\rm req}$ 求成本--生态可行前沿。
目标比例是政策情景选择，\textbf{不是题面给定权重}。敏感性比较保持同一个绝对生态目标，
防止通过降低保护要求制造“更便宜”的假象。整数跳变和求解容差下的有限点列不等于光滑、完整的帕累托曲线。

给定灌溉策略后，任务一的非线性动态先生成每单位面积的逐龄系数；
在空间可加且不共享土壤状态的假设下，目标和资源占用关于面积块、光伏块均为线性，
因此得到混合整数线性模型。最优性证书\textbf{只对策略库与面积网格内的模型有效}，
不外推为全体连续水量路径的全局最优。必须记录可行性、最优界和间隙；
若时间限制内只有可行解，不能写成已证明的精确最优。最大生态参考也受相同界限约束。
用更细面积网格与恒定定额子库对照，揭示离散和策略空间的影响。

\subsection{可执行判据与适用边界}
候选光伏建设年 $s$ 的独立财务筛查为
\begin{equation}\label{eq:threshold}
p_e\ge p_e^{\rm BE}(s)=
\frac{c_p+\sum_{t=s+1}^{T-1}\delta^{t+1-s}(c_{pm}+c_ww_c)}
{\sum_{t=s+1}^{T-1}\delta^{t+1-s}e_0(1-d_p)^{t-s-1}}.
\end{equation}
这是独立项目收支平衡条件，\textbf{不是系统最优的充分条件}。
立项还要检查全期预算、年度预算、每月水量、土地和劳动余量，并在固定生态目标下
比较“包含/不包含该项”的全期最优净成本及求解界。资源紧张时有机会成本，
不能仅用正独立NPV决定实施。没有候选解时应报告不可行或尚无可行解，不补造数值。

模型不计融资、所得税、残值及期末后现金流；不能据此断言工程第 $T$ 年后可以停止养护。
有限期截断、初始植被、储水、策略库、未核实的工程价格及取水许可均限制结论外推。
任务一的长期情景也不能在任务二里突然变成精准预测。

\begin{table}[H]\centering\small
\begin{tabularx}{\textwidth}{lXl}
\toprule
参数 & 含义与证据类别 & 单位\\\midrule
$c_b,c_m$ & 生态建设（本地异类项目参考迁移）、运维（待定价） & 元/hm$^2$；元/(hm$^2$年)\\
$c_w$ & 综合供水成本，未核实本地合同 & 元/m$^3$\\
$c_p,c_{pm}$ & 光伏资本及年运维，工程情景 & 元/kW；元/(kW年)\\
$e_0,p_e,d_p$ & 净上网比产出、电价、衰减；衰减有历史文献参考 & kWh/(kW年)；元/kWh；年$^{-1}$\\
$B_{\rm tot},B_t$ & 总预算与年度额度，规划假设 & 元\\
$Q_t^{\max},Q_{t,m}^{\max}$ & 年/月外部供水上限，需取水批复 & m$^3$\\
$A^{\max},a_t^{\max}$ & 合法用地与年建设能力，规划假设 & hm$^2$\\
$L_t^{\max},\ell_\cdot$ & 可用劳动、单位建设或运维工效，待工时调查 & 人日及相应单位工效\\
\bottomrule
\end{tabularx}
\caption{任务二参数类别；未核实的数值不得列为观测}\label{tab:t2par}
\end{table}

检验至少涵盖光伏电价、建设费、投资比例下界、存量运维、供水、预算与劳动限制，
并对初始植被、有效储水容量和面积离散作结构性敏感性分析。建模可实现并不等于
工程参数已齐备；在本地合同、用地与水量批复缺失时，只能报告条件规划情景。
